\documentclass[aspectratio=169,10pt,spanish]{beamer}

\input{../../../_preambulo_beamer.tex}
\input{../../../_comandos_beamer.tex}

\selectlanguage{spanish}

\title{MA1001B - Modelación Estadística para la Toma de Decisiones}
\subtitle{Lección 15: Aplicaciones binomiales y riesgo de negocio}
\author{}
\institute{Tecnol\'ogico de Monterrey}
\date{}

\begin{document}

\begin{frame}[plain]
  \vspace{1.2cm}
  {\Large\bfseries MA1001B - Modelación Estadística para la Toma de Decisiones\par}
  \vspace{0.7cm}
  {\large Lección 15: Aplicaciones binomiales y riesgo de negocio\par}
  \vspace{0.7cm}
  {\color{TecRojo}\rule{\textwidth}{1.2pt}}
  \vspace{0.8cm}

  Facultad MA1001B\\[0.3cm]
  Periodo\\[0.3cm]
  Tecnol\'ogico de Monterrey
\end{frame}

\begin{frame}{Dos políticas, una media}
  \begin{alertblock}{Pregunta guía}
    Si dos políticas binomiales de inspección tienen ambas $np=2$, ¿tienen
    la misma chance de muchos defectuosos?
  \end{alertblock}

  \vspace{0.6em}
  Al final de esta lección, usted puede:
  \begin{itemize}
    \item enunciar la Política A ($n=50$, $p=0.04$) y la Política B ($n=20$, $p=0.10$);
    \item calcular $P(X\ge 3)$ para cada una;
    \item comparar $P(X=0)$ y $P(X\ge 8)$;
    \item explicar por qué igualar $np$ no iguala el riesgo de cola.
  \end{itemize}

  \vfill
  \scriptsize Todas las políticas usadas hoy son sintéticas. No se usan datos reales de empresa.
\end{frame}

\begin{frame}{Misma $np$, distintos $n$ y $p$}
  \small
  \begin{center}
    \begin{tabular}{lccc}
      \toprule
      \textbf{Política} & $n$ & $p$ & $np$ \\
      \midrule
      A & 50 & 0.04 & 2 \\
      B & 20 & 0.10 & 2 \\
      \bottomrule
    \end{tabular}
  \end{center}

  \vspace{0.5em}
  Varianzas: $\Var(X_A)=1.920000$ y $\Var(X_B)=1.800000$. La Política A tiene
  más espacio para conteos extremos porque $n$ es más grande.
\end{frame}

\begin{frame}[fragile]{Paso 1: Dos políticas}
  \begin{columns}[T]
    \begin{column}{0.42\textwidth}
      \small
      \begin{lstlisting}
n_a, p_a = 50, 0.04
n_b, p_b = 20, 0.10
mean_a = n_a * p_a
mean_b = n_b * p_b
      \end{lstlisting}

      \begin{block}{Salida verificada}
        $np_A=2.000000$\\
        $np_B=2.000000$\\
        $\Var_A=1.920000$\\
        $\Var_B=1.800000$
      \end{block}
    \end{column}
    \begin{column}{0.58\textwidth}
      \centering
      \includegraphics[width=\textwidth]{../figuras/riesgo_binomial_01_dos_politicas.png}
      \vspace{0.2em}
      \scriptsize Ingredientes de las políticas del Paso 1
    \end{column}
  \end{columns}
\end{frame}

\begin{frame}{$P(X\ge 3)$ no es toda la historia del riesgo}
  \small
  \[
    P(X_A\ge 3)=0.323286,\qquad P(X_B\ge 3)=0.323073.
  \]
  Esos valores de cola cercana casi coinciden. La cola lejana no:
  \[
    P(X_A\ge 8)=0.000781,\qquad P(X_B\ge 8)=0.000416.
  \]
  La razón es $1.880028$. La Política A también tiene la mayor chance de
  muestra limpia $P(X=0)=0.129886$ frente a $0.121577$.
\end{frame}

\begin{frame}[fragile]{Paso 2: Probabilidades de cola}
  \begin{columns}[T]
    \begin{column}{0.40\textwidth}
      \small
      \begin{lstlisting}
p_ge3 = stats.binom.sf(2, n, p)
p_ge8 = stats.binom.sf(7, n, p)
      \end{lstlisting}

      \begin{block}{Salida verificada}
        $P(X\ge 3)\approx 0.323$ ambas\\
        $P_A(X\ge 8)=0.000781$\\
        $P_B(X\ge 8)=0.000416$
      \end{block}
    \end{column}
    \begin{column}{0.60\textwidth}
      \centering
      \includegraphics[width=\textwidth]{../figuras/riesgo_binomial_02_probabilidades_cola.png}
      \vspace{0.2em}
      \scriptsize $P(X=0)$ y $P(X\ge 8)$ del Paso 2
    \end{column}
  \end{columns}
\end{frame}

\begin{frame}{Paso 3: Igualar $np$ oculta una cola lejana más pesada}
  \begin{columns}[T]
    \begin{column}{0.42\textwidth}
      \small
      Superponga las dos FMP. La región $x\ge 8$ es casi el doble de probable
      bajo la Política A.

      \vspace{0.4em}
      \textbf{Límite:} la misma media $np$ puede ocultar distintas
      probabilidades de cola. La Lección 16 extrae sin reemplazo, de modo que
      los ensayos ya no son independientes.
    \end{column}
    \begin{column}{0.58\textwidth}
      \centering
      \includegraphics[width=\textwidth]{../figuras/riesgo_binomial_03_misma_media_colas_distintas.png}
      \vspace{0.2em}
      \scriptsize Superposición de FMP del Paso 3
    \end{column}
  \end{columns}
\end{frame}

\begin{frame}{Verificación de conceptos}
  \begin{enumerate}
    \item ¿Por qué ambas políticas tienen media $2$?
    \item Calcule $P(X\ge 3)$ para la Política A a partir de la salida verificada.
    \item ¿Por qué $P(X\ge 3)$ puede coincidir cuando $P(X\ge 8)$ no coincide?
    \item ¿Cuál política tiene la cola lejana más pesada?
  \end{enumerate}

  \vspace{0.6em}
  \begin{block}{Conexión con la Lección 16}
    La sesión puede continuar directamente con el modelo hipergeométrico para
    muestreo sin reemplazo.
  \end{block}

  \vfill
  \scriptsize Fuente: OpenStax, \textit{Introductory Business Statistics 2e},
  Sección 4.2. CC BY-NC-SA.
\end{frame}

\appendix



\end{document}
